\begin{theorem}[DP34: optimal constant-one affine bounds on every Honda dual]\label{DP34}
In the actual-dual coordinates of DP32, for every $h\ge2$, every $i\in\{1,\ldots,h-1\}$ and every integral $f$,
\[
 D_i^nf\in p^{b_i(n)}R[[\mathbf X]],\qquad
 b_i(n)=\max\{0,\ceil{c_i n-1}\},\qquad c_i=\frac{p^i}{p^h-1}.
\]
Both the slope and the constant $1$ are optimal. With $q=p^h$ and $d=h-1$, the same monomial
\[
 f_*(\mathbf X)=\left(\prod_{j<d}X_j^{p-1}\right)X_d^{p^2-1}
\]
satisfies $v_p(D_i^{q-1}f_*(0))=p^i-1$ for every $i$. Mixed derivatives have the sum $\sum_i b_i(a_i)$ as a lower bound. The displayed ideal containments persist under integral base change.
\end{theorem}
\begin{proof}
Work first over $\Z_p$ in the model already identified by DP32. Integral translation coefficients and the first nonlinear terms of its logarithm give, exactly as in DP33,
\begin{equation}\label{eq:all-dual-low}
 D_i^p=aD_{i+1}+pT_i\quad(i<d),\qquad
 D_d^{p^2}=p^pbD_1+p^{p+1}T_d,\qquad D_d^p\in p\operatorname{End}_{\rm int}.
\end{equation}
Here $a=-(p-1)!$, $b=-(p^2)!/p^{p+1}$ are units; all remainder operators preserve the integral function ring and commute. The last two identities also hold when $d=1$. Descending induction, raising the first identity to $p^{h-i-1}$, shows that $D_i^{R_i}\in p\operatorname{End}_{\rm int}$ for $R_i=p^{h-i}$. Hence $D_i^n$ has the preliminary divisibility $\floor{n/R_i}$.

Fix $i<d$, put $B=p^i$, and start with $D_i^q$. Successively substitute the first identity of \eqref{eq:all-dual-low} to move its leading term through directions $j=i,i+1,\ldots,d-1$. At the step $j=i+t$, this identity is raised to $N=p^{h-t-1}$, and $N/R_{j+1}=B$. Its $m$-th nonleading term has valuation at least
\[
 v_p\binom Nm+m+\floor{(N-m)/R_{j+1}}
 =B+v_p\binom Nm+m-\ceil{m/R_{j+1}}.
\]
For $1\le m<N$, the binomial coefficient is divisible by $p$ because $N$ is a $p$-power, and $m-\ceil{m/R_{j+1}}\ge0$. For $m=N$ the bound is $N\ge p^{i+1}\ge B+1$. Thus every discarded term belongs to $p^{B+1}\operatorname{End}_{\rm int}$. The leading term at the end is a unit times $D_d^{p^{i+1}}$. Raise the second identity of \eqref{eq:all-dual-low} to $p^{i-1}$ and factor out $p^B$. Modulo $p$, what remains is a unit times $D_1^{p^{i-1}}$, which the first identities modulo $p$ identify with a unit times $D_i$. Therefore
\begin{equation}\label{eq:all-dual-q}
 D_i^q=p^{p^i}u_iD_i+p^{p^i+1}U_i
\end{equation}
with a unit $u_i$ and an integral operator $U_i$. For $i=d$, use the same last step directly by raising the second identity to $p^{d-1}$. This includes $d=1$.

Now $c_iR_i=q/(q-1)$ and $c_i(q-1)=p^i$. For $0\le n<q$, the increasing block maxima give $\floor{n/R_i}\ge c_i n-1$, with equality in the real difference at $n=q-1$. For $n\ge q$, applying \eqref{eq:all-dual-q} to $D_i^{n-q}f$ bounds its first term by $c_i n-1$ and its second by $c_i n-1+(1-c_i)\ge c_i n-1$. This proves the affine bound by induction. All coefficients and operators were over $\Z_p$, so integer valuations give the asserted ideal containment there. Base change preserves that containment; no fractional valuation is rounded on a ramified base. Composition proves the mixed bound.

It remains to prove the sharp common witness, including exclusion of possible lower-degree terms. The explicit logarithm and its inverse satisfy the Teichm\"uller equivariance
\[
 \mathcal L_j(\zeta^pX_1,\zeta^{p^2}X_2,\ldots,\zeta^{p^d}X_d)
 =\zeta^{p^j}\mathcal L_j(\mathbf X),\qquad\zeta^{q-1}=1.
\]
For the ordinary arrows this is $p\cdot p^{j-1}=p^j$; for the exceptional arrow it is $p^2p^d\equiv p\pmod{q-1}$. Thus a nonzero degree $n$ in $\mathcal E_j(te_i)$ requires $np^i\equiv p^j\pmod{q-1}$. Its smallest possible positive degree is
\[
 m_{ij}=p^{j-i}\ (j\ge i),\qquad m_{ij}=p^{h+j-i}\ (j<i).
\]
DP32's scaled inverse gives a nonzero coefficient at this first path degree, with exact valuation $r_j-m_{ij}r_i$. Equivariance excludes all smaller degrees. If $e_j=p-1$ for $j<d$ and $e_d=p^2-1$, then
\[
 \sum_j e_jm_{ij}=q-1,\qquad \sum_j e_jr_j=d.
\]
The first equality telescopes: the ranges $i\le j<d$, $j=d$, and $j<i$ contribute $p^{d-i}-1$, $p^{h-i+1}-p^{d-i}$, and $q-p^{h-i+1}$ respectively. For the second, use $\sum e_j=(p-1)(h+p-1)$ and $\sum e_jp^j=p(q-1)$. The unique lowest-degree term of $f_*(\mathcal E(te_i))$ therefore has degree $q-1$ and valuation $d-(q-1)r_i$. Finally
\[
 v_p((q-1)!)=\frac{q-1}{p-1}-h
\]
makes the derivative valuation $c_i(q-1)-h+d=p^i-1$. This proves the optimal constant in every direction. Slope optimality is already the exact sequence in DP32. Every step applies at $p=2$ and $h=2$.
\end{proof}
